
{% QClaudePyMisc137
\needspace{17\baselineskip}
\item \rtask \ponto{\pt}%
Em Python~3, a função embutida \lstinline[style=Python]|tuple()|
pode converter qualquer iterável em uma tupla. 

Sobre a conversão
de \lstinline[style=Python]|set| para \lstinline[style=Python]|tuple|,
assinale a alternativa CORRETA:

\begin{answerlist}[label={\texttt{\Alph*}.},leftmargin=*]
    \ti A conversão preserva a ordem de inserção dos
        elementos no conjunto original.

    \di A conversão é possível e válida, porém a
        {ordem dos elementos} na tupla resultante
        pode não corresponder à ordem de inserção no
        conjunto, pois \lstinline[style=Python]|set| é uma
        coleção não ordenada.

    \ti \lstinline[style=Python]|tuple(set)| levanta
        \lstinline[style=Python]|TypeError| pois
        \lstinline[style=Python]|set| não é iterável.

    \ti A conversão sempre produz a tupla em ordem crescente
        dos elementos.

    \ti A conversão é impossível porque \lstinline[style=Python]|set|
        pode conter elementos não \textit{hashable}.
\end{answerlist}
  % ----------------------------------------------------------------------
  % ANÁLISE TÉCNICA
  % ----------------------------------------------------------------------
  %
  % tuple(set_obj):
  %   → Válido: set é iterável, então tuple() o percorre.
  %   → A ordem na tupla resultante é a ordem de iteração do set.
  %   → Set NÃO garante ordem; depende dos hashes e da implementação.
  %   → A ordem pode variar entre execuções (PYTHONHASHSEED).
  %
  % Exemplo:
  %   s = {3, 1, 2}
  %   tuple(s) → pode ser (1, 2, 3) ou (3, 1, 2) ou outra ordem.
  %
  % Diferença importante:
  %   tuple(sorted({3,1,2})) → (1,2,3) → ordem garantida (sorted).
  %   tuple({3,1,2}) → ordem não garantida.
  %
  % Por que A está errada: set não preserva ordem de INSERÇÃO.
  % Por que D está errada: tuple(set) não é necessariamente crescente.
  %
  % ----------------------------------------------------------------------
  % Análise das alternativas:
  % ----------------------------------------------------------------------
  %
  % A) Incorreta. set não preserva ordem de inserção.
  % B) Correta.  Conversão válida; ordem não garantida (set é desordenado).
  % C) Incorreta. set é iterável; tuple(set) é válido.
  % D) Incorreta. A ordem não é necessariamente crescente.
  % E) Incorreta. set armazena apenas elementos hashable; conversão é válida.
  %
  % ----------------------------------------------------------------------
  % CONCLUSÃO
  % ----------------------------------------------------------------------
  %
  % Alternativa correta: B
  % ----------------------------------------------------------------------
}
